\section{Protocole d’évaluation}
\label{sec:evaluation}

Rien de ce qui précède ne montre que la carte aide à lire. Le protocole qui suit est prévu ; aucun résultat n’est encore disponible.

\paragraph{Qualité de la reconstruction.}
Une dizaine de textes couvrant les trois genres du sujet, articles scientifiques, textes philosophiques et essais, seront annotés indépendamment par deux personnes : étapes principales, relations et phrases d’appui. L’accord entre annotateurs sera mesuré avant toute comparaison avec le modèle, car il peut être faible pour les affirmations \parencite[631]{ref07} et les dépendances à longue distance échappent aux annotateurs eux-mêmes \parencite[43]{ref08}. On mesurera ensuite, pour chaque carte produite, la correspondance des étapes avec l’annotation, la justesse du type et du sens des relations, et la part des ancrages qui soutiennent réellement leur étape. Cette dernière mesure servira aussi à calibrer les statuts : parmi les éléments marqués \enquote{vérifié}, quelle part l’est aussi pour les annotateurs ?

\paragraph{Utilité pour la lecture.}
Une étude intra-sujets comparera trois conditions : le texte seul, le texte accompagné d’un résumé produit par le même modèle, et le texte accompagné de la carte. Les tâches consisteront à retrouver la phrase qui appuie une affirmation donnée, à repérer une conclusion insuffisamment soutenue et à répondre à des questions sur la structure du raisonnement. On mesurera l’exactitude, le temps et la confiance déclarée, rapportée à l’exactitude. Une confiance élevée sur des réponses fausses serait un mauvais signe, compte tenu de ce qu’a observé Story Ribbons \parencite[8]{ref01}.

\paragraph{Menaces à la validité.}
Trois points appellent une attention particulière. La familiarité avec les textes : Story Ribbons laissait chaque participant choisir une œuvre qu’il connaissait \parencite[7]{ref01}, ce qui facilite le jugement des sorties du modèle mais ne représente pas la lecture d’un texte nouveau. L’effet de nouveauté et la brièveté des séances, que les auteurs de Sensecape signalent comme limites de leur propre étude \parencite[13]{ref03}. Enfin, la variabilité des sorties du modèle d’une exécution à l’autre, qui rend difficile toute évaluation quantitative globale \parencite[17]{ref04} : les cartes utilisées dans l’étude seront figées à l’avance, et leur variabilité mesurée séparément.
